%----------------------------------------------------------------------------------------
%	Metropolia Thesis LaTeX Template
%----------------------------------------------------------------------------------------
% License:
% This work is licensed under the Creative Commons Attribution 4.0 International License.
% To view a copy of this license, visit http://creativecommons.org/licenses/by/4.0/.
%
% However, this license apply to this template. As a template, it is supposed to be
% modified for your own needs (with your thesis content). For this reason, if you use
% this project as a template and not specifically distribute it as part of a another
% package/program, we grant the extra permission to freely copy and modify these files as
% you see fit and even to delete this copyright notice.
% In short, you are free to publish your thesis under whatever license you wish, even
% keep the all rights reserved to you.
%
% Authors:
% Panu Leppäniemi, Patrik Luoto, Mikaa Oni and Patrick Ausderau
%
% Credits:
% Panu Leppäniemi: abstract, def, cleaning,...
% Patrik Luoto: title page, abstract in Finnish, abbreviation, math,...
% Mikaa Oni: switch to biber biblatex
% Patrick Ausderau: initial version, style, table of content, bibliography, figure,
%                   appendix, table, source code listing...
%
% Please:
% If you find mistakes, improve this template and alike, please contribute by sharing
% your improvements and/or send us your feedback there:
% https://github.com/panunu/metropolia-thesis-latex
% And of course, if you improve it, add yourself as an author.
%
% Compiler:
% Use XeLaTeX as a compiler.

%----------------------------------------------------------------------------------------
%	THESIS INFO
%----------------------------------------------------------------------------------------

% All general information (main language, title, author (you), degree programme, major
% option, etc.)
% Edit the file chapters/0info.tex to change these information
\documentclass[12pt,a4paper,oneside,article]{memoir}%Do not touch this first line ;)

% Global information (title of your thesis, your name, degree programme, major, etc.)

\def\bilingual{no}%For Finnish students, you must have 2 abstracts, one in English and one in your native language (Finnish or Swedish), so "yes", your thesis is bilingual.
%\def\bilingual{no}%For international student writing in English, only one language and one abstract.

\def\thesislang{english} %change this depending on the main language of the thesis.
%\def\thesislang{english} % "english" is the only other supported language currently. If someone has the swedish, please contribute!

%\def\secondlang{english} %if the main language is Finnish (or Swedish), you must have 2 abstracts (one in Finnish (or Swedish) and one in English)
%\def\secondlang{finnish}
%If the main language is English and that you are native Finnish (or Swedish) speaker, you must have also abstract in your native language on top of the English one.

\author{Bhabani Sankar Das} %your first name and last name

%\def\alaotsikko{Alaotsikko/Subtitle} %DISABLED, seems not to be an option with the 2018 template. If you really need it, uncomment and modify style/title.tex accordingly.

%License
%When publishing your thesis to theseus.fi, you can keep all rights reserved to you,
%or use one of the Creative Commons https://creativecommons.org/licenses/?lang=en
%This attribute will set the license in the metadata of the generated pdf.
%possible options (case sensitive):
%all (keep all rights reserved to yourself)
%by (Attribution)
%by-sa (Attribution-ShareAlike)
%by-nd (Attribution-NoDerivs)
%by-nc (Attribution-NonCommercial)
%by-nc-sa (Attribution-NonCommercial-ShareAlike)
%by-nc-nd (Attribution-NonCommercial-NoDerivs)
%Note that theseus do not accept dedication to public domain CC0
\def\thesiscopy{all}

%Finnish section, for title/abstract
\def\otsikko{Opinnäytetyön otsikko}
\def\tutkinto{Tutkinto (esim. Insinööri (AMK))} % change to your needs, e.g. "YAMK", etc.
\def\kohjelma{Koulutusohjelma (esim. Tieto\textendash ja viestintätekniikka)}
\def\suuntautumis{Ammatillinen pääaine (esim. Mobile Solutions)}
\def\thesisfi{Insinöörityö}%was Opinnäytetyö
\def\ohjaajat{
Titteli Etunimi Sukunimi\newline
Titteli Etunimi Sukunimi
}
\def\tiivistelma{
Tämä on tiivistelmän ensimmäinen kappale. Tiivistelmän kappaleet loppuvat komentoon newline, jotta saadaan yksi tyhjä rivi aikaiseksi. \newline

Tämä on tiivistlemän toinen kappale.
}
\def\avainsanat{avainsana, avainsana}
\def\aihe{Lyhyt kuvaus opinnäytetyöstä. Max 255 merkkiä.}%for the pdf metadata/properties. If not used, empty it and also the \def\subject.

%English section, for title/abstract
\title{Improvement of Embedded System Testing by Adding Fuzzing Technique- A Case Study}
\def\metropoliadegree{Master of Engineering} % change to your needs, e.g. "master", etc.
\def\metropoliadegreeprogramme{Information Technology}
\def\metropoliaspecialisation{Networking and Services }
\def\thesisen{Master’s Thesis} % change to your need, e.g. master's
\def\metropoliainstructors{
First name Last name, Title (e.g.: Project Manager)\newline
First name Last name, Title (Sami Sainio: Principal Lecturer)
}
\def\abstract{
System testing of any software product is very important to maintain good quality, security, and
integrity. Embedded systems which bind the software execution tightly between the hardware
architecture and software in a tight loop has grown at a steady pace over the years.
As the embedded systems are complex, it is recommended to apply software testing throughout the
development process.

The manual, traditional and automated testing of embedded systems often follow the same static
approaches to uncover any vulnerabilities. Fuzzing or fuzz testing is a dynamically automated
testing technique which subjects a system to a stream of input data by modifying the valid seeds.
Any program which accepts untrusted and malformed data is ideal for applying fuzzing. Fuzzing has
been used to exploit any edge cases by generating huge number of input data. The errors are
not essentially related to functional requirements.

This thesis explores the possibility of the adoption of fuzzing in the embedded software development
lifecycles to discover any faults. The aim is to do a literature review of different fuzzing
techniques, methods, and the tools available. For the implementation of the case study the project
``OSS-Fuzz'' was chosen.

The output and results of the thesis are that fuzzing is a popular way to test the correctness and
security of embedded systems and can be implemented for the overall improvement in the embedded system testing.
Found out that, the Fuzz testing must be deployed continuously to find the vulnerabilities before the software get shipped.
Fuzzing helps in findings bugs irrespective of the project requirements and not a replacement to the functional
testing, but a must necessary add-on.


}
\def\metropoliakeywords{Fuzzing, Embedded system, crypto, software testing}
\def\subject{short description of the thesis. Max 255 characters.}%for the pdf metadata/properties. If not used, empty it and also the \def\aihe.


%----------------------------------------------------------------------------------------
%	GLOBAL STYLES
%----------------------------------------------------------------------------------------

% If you need extra package, etc. modify the style/style.tex file.
% If you are using Windows OS, you will need to change default font to Arial in that
% style/style.tex file (or install Liberation Sans font to your system).
% If you are using MacOS or linux, make sure you have Liberation Sans font installed.
% Global style. Normally should not be edited.
% If you use windows OS, eventually change \setmainfont to Arial
% Check around commit https://github.com/panunu/metropolia-thesis-latex/commit/a0c15ac77bab1a52c59c517a18080938e57bf5ef
% to see how the font files were manually added (after downloading them: https://pagure.io/liberation-fonts/ )

%\usepackage[l2tabu, orthodox]{nag}%check for obsolete packages (with outdated nag package?)

\usepackage[\secondlang,\thesislang]{babel}% finnish english swedish
\usepackage{iflang}
\usepackage{amsmath}
\usepackage{amsfonts}
\usepackage{amssymb}
\usepackage{fontspec}
\usepackage{tocloft}
\usepackage{titlesec}
\usepackage{mathtools}
\usepackage[amssymb]{SIunits}
\usepackage[version=3]{mhchem}
\usepackage{tikz} % mindmaps, flowcharts, piecharts, examples at http://www.texample.net/tikz/examples/
\usetikzlibrary{shapes.geometric, arrows}

%condition for adding or not space in TOC
\usepackage{etoolbox}
%for compact list
\usepackage{enumitem}
%for block comment
\usepackage{verbatim}
%for "easier" references
%\usepackage{varioref}
%forcing single line spacing in bibliography
\DisemulatePackage{setspace}
\usepackage{setspace}
%including figure (image)
\usepackage{graphicx}
%change the numbering for figure
\usepackage{chngcntr}
%strike trough
\usepackage{ulem}
%euro symbol
\usepackage{eurosym}
%try to count
\usepackage{totcount}
%insert source code
%\usepackage{listings}
%require -8bit -shell-escape in the xelatex compile command
%if compiling locally, consider options cachedir=minted,outputdir=~/.tex
\usepackage[newfloat]{minted}
\setminted{tabsize=2,linenos,breaklines,breaksymbolleft={\quad},baselinestretch=1}
\setmintedinline{breaklines}
\usepackage[singlelinecheck=false]{caption}
\usepackage{color}
%force the width of a table instead of column
\usepackage{tabularx}
\usepackage{booktabs} %why not booktabs? :3

\usepackage{csquotes}% avoid warning with babel

\usepackage{float} % For forced figure location with modifier H (\begin{figure}[H])

% citep-macro for reference with period inside square brackets [1.]
\newcommand{\citep}[1]{
 \renewcommand\citeright{.]}
 \cite{#1}
 \renewcommand\citeright{]}
}

%set date format to D.M.YYYY
\def\pvm{\the\day.\the\month.\the\year}
%set date format to D Month YYYY
\usepackage[en,useregional=false]{datetime2}
\DTMsetup{datesep=.}
\DTMnewdatestyle{dMonthyyyy}
{%
  \renewcommand{\DTMdisplaydate}[4]{%
    \number##3 % day
    ~% separator
    \DTMenglishmonthname{##2}% month name
    ~% separator
    \number##1% year
  }%
  \renewcommand{\DTMDisplaydate}[4]{%
    \DTMdisplaydate{##4}{##3}{##2}{##1}%
  }%
}
\DTMsetdatestyle{dMonthyyyy}
\date{\today}

\newcommand\tn[1]{\textnormal{#1}} %use \tn instead of \textnormal
\newcommand\reaction[1]{\begin{equation}\ce{#1}\end{equation}} %\reaction{} for chemical reactions

%NORMAL TEXT
%all text, title, etc. in the same font: Arial
%NOTE: fontname is case-sensitive
\setmainfont[Scale=0.98]{Liberation Sans}
%line space
\linespread{1.46}
\AtBeginEnvironment{tabular}{\singlespacing}
%\doublespacing
%margin
%geometry moved after hyperref
\setlength{\parindent}{0pt} %first line of paragraph not indented
\setlength{\parskip}{16.5pt} %one empty line to separate paragraph
%list with small line space separation
\tightlists

%IMAGE - FIGURE
%the figures should be placed in the "illustration" folder
\graphicspath{{illustration/}}
%figure number without chapter (1.1, 1.2, 2.1) to (1, 2, 3)
\counterwithout{figure}{chapter}
%border around images
\setlength\fboxsep{0pt}
\setlength\fboxrule{0.5pt}
%space after figure caption (and other float elements)
\setlength{\belowcaptionskip}{-7pt}

%TABLE
\counterwithout{table}{chapter}

%SOURCE CODE
\newenvironment{code}{\captionsetup{type=listing}}{}
\IfLanguageName{finnish}{\SetupFloatingEnvironment{listing}{name=Koodiesimerkki}}{}%was Listaus
%\counterwithout{lstlisting}{chapter}
%moved after begin document, otherwise does not compile

%% set this format as the default for lstlisting
%\DeclareCaptionFormat{empty}{}
%\captionsetup[lstlisting]{format=empty}

%TOC
%change toc title
\IfLanguageName{finnish}{\addto{\captionsfinnish}{\renewcommand*{\contentsname}{Sisällys}}}{}
%remove dots
\renewcommand*{\cftdotsep}{\cftnodots}
%chapter title and page number not in bold
\renewcommand{\cftchapterfont}{\normalfont}
\renewcommand{\cftchapterpagefont}{\normalfont}
%sub section in toc
\setcounter{tocdepth}{2}
%subsection numbered
\setcounter{secnumdepth}{2}
\renewcommand{\tocheadstart}{\vspace*{-33pt}}
\renewcommand{\printtoctitle}[1]{\fontsize{13.5pt}{13.5pt}\bfseries #1}
%\renewcommand{\aftertoctitle}{\vspace*{-22pt}\afterchaptertitle}
%spacing after a chapter in toc
\preto\section{%
  \ifnum\value{section}=0\addtocontents{toc}{\vskip11pt}\fi
}
%spacing after a section in toc
\renewcommand{\cftsectionaftersnumb}{\vspace*{-3pt}}
%spacing after a subsection in toc
\renewcommand{\cftsubsectionaftersnumb}{\vspace*{-1pt}}
%appendix in toc with "Appendix " + num
\IfLanguageName{finnish}{
  \renewcommand*{\cftappendixname}{Liite\space}
  \renewcommand{\appendixtocname}{Liitteet}
}{\renewcommand*{\cftappendixname}{Appendix\space}}
%appendix header
\IfLanguageName{finnish}{\def\appname{Liite\space}}{\def\appname{Appendix\space}}

%TITLES
%chapter title
%\clearforchapter{\clearpage}
\titleformat{\chapter}
{\fontsize{14pt}{14pt}\bfseries\linespread{1}}%\clearpage
{\thechapter}{.5cm}{}
\titlespacing*{\chapter}{0pt}{-.42cm}{.5pt}
\titleformat{\section}
{\fontsize{13.5pt}{13.5pt}\normalfont}
{\thesection}{.5cm}{}
\titlespacing*{\section}{0pt}{9pt}{0pt}
\titleformat{\subsection}
{\fontsize{12.7pt}{12.7pt}\normalfont}
{\thesubsection}{.5cm}{}
\titlespacing*{\subsection}{0pt}{11pt}{0pt}


%QUOTE
\renewenvironment{quote}
  {\list{}{\rightmargin=0pt\leftmargin=1cm\topsep=-10pt}%
  \item\relax\fontsize{10pt}{10pt}\singlespacing}
  {\endlist}

%BIBLIOGRAPHY
%bibliography title to be "references"
%IF THE TITLE DON'T GET RENAMED PROPERLY, move that line after the \begin{document}
\IfLanguageName{finnish}{\addto{\captionsfinnish}{\renewcommand*{\bibname}{Lähteet}}}{\renewcommand\bibname{References}}
\makeatletter %reference list option change
\renewcommand\@biblabel[1]{#1\hspace{1cm}} %from [1] to 1 with 1cm gap
\makeatother %
\setlength{\bibitemsep}{11pt}

\usepackage[backend=biber,bibencoding=utf8,%
citetracker=true,%
isbn=true,%
doi=true,%
url=true,%
usetranslator=true,%
style=numeric-comp,%
citestyle=numeric,%
terseinits=false,%
giveninits=true,%
sorting=none%
]{biblatex}

% set right format
%\DeclareNameAlias{sortname}{last-first} % deprecated
\DeclareNameAlias{default}{family-given}
\DeclareFieldFormat{labelnumberwidth}{#1} % remove () from label number
\DeclareFieldFormat{title}{"#1"} % quotes to the title
\DeclareFieldFormat{journaltitle}{#1} % remove underline
\DeclareFieldFormat*{url}{\textless\url{#1}\textgreater} % you can modify how to url looks here
\DeclareFieldFormat{urldate}{\addcomma\space\bibstring{urlseen}\space#1} % remove () from date
%try set translation to biblio
\DefineBibliographyStrings{english}{%
    urlfrom = {available at},%
    urlseen = {visited on},%
    fromenglish = {from English},%
    fromfinnish = {from Finnish},%
    fromgerman = {from German},%
    fromjapanese = {from Japanese},%
}
\DefineBibliographyStrings{finnish}{%
  %  urlfrom={Linkki: },%
    urlfrom = {},%
    urlseen = {katsottu},%
    fromjapanese = {japanista},%
    fromenglish = {englannista},%
    fromfinnish = {suomesta},%
    fromgerman = {saksasta},%
}
{
  %new cite command: "Vancouver Short"
  \DeclareCiteCommand{\citeVS}
    {\usebibmacro{prenote}}
    {\usebibmacro{author}, \usebibmacro{title}}
    {\multicitedelim}
    {\usebibmacro{postnote}}

  % new cite command: "Vancouver Short Collection" - necessary when referencing whole collections.
  \DeclareCiteCommand{\citeVSc}
    {\usebibmacro{prenote}}
    {\usebibmacro{editor}, \usebibmacro{title}}
    {\multicitedelim}
    {\usebibmacro{postnote}}
}

\addbibresource{biblio.bib}% for biblatex you need out \printbibliography too

%count the appendices (since the chapter counter is reset after \appendix).
%! require to complie 2 times
\regtotcounter{chapter}

% metadata (title, author, lang,...) for accessibility, etc.
\usepackage{hyperref}
\usepackage{hyperxmp}
\def\isolang{\IfLanguageName{finnish}{fi}{en}} %iso code (based on main language)
%TODO bug: pdfcopyright don't go, neither the alt lang
%works with simple example, so probably conflict with some package?
\hypersetup{%
  pdfdisplaydoctitle,
  %breaklinks=true,%searching for overfull warnings
  pdfencoding=auto,
  bookmarksdepth=subsection,
  unicode=true,
  keeppdfinfo=true,
  pdflang={\isolang},
  pdfmetalang={\isolang},
  pdftitle={\IfLanguageName{finnish}{\otsikko}{\thetitle}},
  pdfkeywords={\metropoliakeywords},
  pdfcopyright={Copyright \textcopyright\ \the\year{}, \theauthor},
}
\XMPLangAlt{\IfLanguageName{finnish}{en}{fi}}{%
  pdftitle={\IfLanguageName{finnish}{\thetitle}{\otsikko}},
  pdfcopyright={Copyright \textcopyright\ \the\year{}, \theauthor},
}
\urlstyle{same}

%moved after hyperred as can cause conflicts
\usepackage{pdfcomment}%try the alt text for graphics
\usepackage{accsupp}
\newcommand{\AltText}[2]{\BeginAccSupp{method=pdfstringdef,unicode,Alt={{#1}}}\pdftooltip{{#2}}{{#1}}\EndAccSupp{}}

\usepackage[top=2.5cm, bottom=3cm, left=4cm, right=2cm, nofoot]{geometry}

\usepackage{pgfplots} %simple plots etc
\pgfplotsset{compat=1.17}
\usepackage{pgfplotstable}

% Abbreviations, acronym and glossary, in case of bug, remove temporary the noredefwarn
\usepackage[acronym,toc,nonumberlist,section=chapter,noredefwarn]{glossaries}%xindy,%toc, ,nomain
\newglossarystyle{mystyle}{%
  \setglossarystyle{list}% base this style on the list style
  \renewcommand*{\glossentry}[2]{%
  \item[\glsentryitem{##1}%
    \glstarget{##1}{\glossentryname{##1}:}]
  \glossentrydesc{##1}\glspostdescription\space ##2}
}
\setglossarystyle{mystyle}

\renewcommand*{\glsclearpage}{}

% Normally, you do not need to modify the title style. It's content comes from the
% chapters/0info.tex file.
% TITLE PAGE
% Normally, you should not edit this file.

\makeatletter
\renewcommand{\maketitle}{
\newgeometry{left=4cm,top=2.7cm}
\thispagestyle{empty}

\AltText{\IfLanguageName{finnish}{Metropolia Ammattikorkeakoulu.}{Metropolia University of Applied Sciences.}}{\includegraphics[width=9.1cm,keepaspectratio]{Metropolia_L_oranssi}}

\IfLanguageName{finnish}{
\vspace*{1.36cm}
}{
\vspace*{1.18cm}
}
\tn{\large\theauthor\\[9.2mm]{\fontsize{27.5pt}{27.5pt}\normalfont\color[RGB]{155,50,35}\IfLanguageName{finnish}{\otsikko}{\thetitle}}}%\\[22pt]\LARGE\alaotsikko\\[1.75cm]}

\vspace*{2.38cm}
{\setstretch{1.512}
\IfLanguageName{finnish}{
  Metropolia Ammattikorkeakoulu\\
  \tutkinto \\
  \kohjelma \\
  \thesisfi\\
  \pvm % D.M.YYYY date format for Finnish.
} {
  Metropolia University of Applied Sciences\\
  \metropoliadegree \\
  \metropoliadegreeprogramme \\
  \thesisen\\
  \@date %D Month YYYY for English
}
}
\clearpage
\restoregeometry
}
\makeatother



%----------------------------------------------------------------------------------------
%	ABBREVIATION AND GLOSSARY
%----------------------------------------------------------------------------------------

% Add/edit all your acronyms, abbreviations, glossary entries, etc. definitions in
% chapters/0abbr.tex file.
% You can have as many as you wish. Only the ones you use in your text (inserted with
% \gls{} command) will print in the Glossary/Lyhenteet.
% Generate the glossary
\makeglossaries

% Acronyms, abbreviations, etc.

\newacronym{html}{HTML}{HyperText Markup Language}
\newacronym{sql}{SQL}{Structured Query Language}
\newacronym{io}{I/O}{Input/Output}
\newacronym{ram}{RAM}{Random Access Memory}
\newacronym{php}{PHP}{Hypertext Preprocessor}
\newacronym{wysiwym}{WYSIWYM}{What You See Is What You Mean}
\newacronym{isbn}{ISBN}{International Standard Book Number}
\newacronym{url}{URL}{Uniform Resource Locator}
\newacronym{doi}{DOI}{Digital Object Identifier}
\newacronym{iot}{IoT}{Internet of Things}
\newacronym{ics}{ICS}{Industrial Control System}
\newacronym{ip}{IP}{Intellectual Property}
\newacronym{rot}{ROT}{Root of Trust}
\newacronym{soc}{SoC}{system-on-chip}
\newacronym{os}{OS}{Operating Systems}
\newacronym{iso}{ISO}{International Organization of Standardization}
\newacronym{dpa}{DPA}{Differential Power Analysis}
\newacronym{sdlc}{SDLC}{Software Development Lifecycle}
\newacronym{stlc}{STLC}{Software Testing Lifecycle}
\newacronym{sut}{SUT}{System Under Test}
\newacronym{smt}{SMT}{Satisfiability Modulo Theory}
\newacronym{ddos}{DDoS}{Distributed Denial of Service}
\newacronym{rtos}{RTOS}{Real Time Operating System}
\newacronym{afl}{AFL}{American Fuzzy Lop}
\newacronym{ci/cd}{CI/CD}{Continuous Integration and Continuous Deployment}
\newacronym{gui}{GUI}{Graphical User Interface}


% Glossary entries

\newglossaryentry{part_key}{
	name={partition key},
	description={a column or set of columns from the same table whose consolidated value decide the partition for a given data}
}
\newglossaryentry{thesis}{
	name=thesis,
	description={a written essay one submitted for a university degree},
	plural=theses
}
\newglossaryentry{latex}
{
	name=\LaTeX{},
	description={Is a mark up language specially suited for scientific documents}
}

\newglossaryentry{DDoS}
{
	name=Distributed Denial of Service,
	description={A malicious attempt to overwhelm a targeted
	system, network, or website by inundating it with a flood of traffic
	from multiple sources}
}

\newglossaryentry{RTOS}
{
	name=Real Time Operating System,
	description={An operating system that is specifically designed to
	ensure real-time response and predictability in embedded systems
	with critical timing requirements}
}
\newglossaryentry{Fuzzers}
{
	name=Fuzzers,
	description={Fuzzers are software testing tools used to find security
	vulnerabilities and software bugs by submitting random or
	semi-random inputs to a program}
}

\newglossaryentry{fuzz_target}
{
	name=Fuzz Target,
	description={A specific function or area of a software application or
	system that is intentionally subjected to fuzz testing}
}
\newglossaryentry{corpus}
{
	name=corpus,
	description={Designed to include set of minimal inputs for fuzzing }
}
\newglossaryentry{fuzzing_engine}
{
	name=Fuzzing Engine,
	description={A tool that finds and executes inputs for a fuzz target}
}

%----------------------------------------------------------------------------------------
%	DOCUMENT STARTS HERE...
%----------------------------------------------------------------------------------------

\begin{document}
\IfLanguageName{finnish}{
}{
  \raggedright%2021 template, align left, no hyphennization for English version
}
\counterwithout{listing}{chapter}

%----------------------------------------------------------------------------------------
%	TITLE PAGE
%----------------------------------------------------------------------------------------

% Title header. You should not edit that file.

%page number always on the top right, clear the "chapter/section" head
\pagestyle{myheadings}
\markright{}
%clear chapter "title" foot page
\makeevenfoot{plain}{}{}{}
\makeoddfoot{plain}{}{}{}



\maketitle
\newpage

%----------------------------------------------------------------------------------------
%	ABSTRACT / Tiivistelmä
%----------------------------------------------------------------------------------------

% If you are international student writing in English, ignore the Finnish abstract.
% If you are Finnish citizen, you must have 2 abstracts, one in Finnish (or Swedish
% depending on your mother tongue) and one in English regardless of the main language of
% your thesis. Normally, you do not need to modify the abstract style. It's content comes
% from the chapters/0info.tex file.
\ifstrequal{\bilingual}{no}{%
    % Abstract in English
% Normally, you should not edit this file. Everything comes from chapter/0info.tex

\pagestyle{empty} %remove page number
\newgeometry{top=2.45cm,left=4cm,right=.5cm}
\begin{otherlanguage}{english}
{\large\textbf{Abstract}}
{\renewcommand{\arraystretch}{1.1}
  \begin{tabular}{@{}p{4.7cm} >{\raggedright\arraybackslash}p{10.8cm}@{}}
  Author: & \makeatletter\@author\makeatother
  \\
  Title: & \makeatletter\@title\makeatother
  \\
  Number of Pages: & \pageref*{LastPage} pages + \total{chapter} appendices
  \\
  Date: & \makeatletter\thedate\makeatother
  \\[6.5mm]
  Degree: & \metropoliadegree
  \\
  Degree Programme: & \metropoliadegreeprogramme
  \\
  Professional Major: & \metropoliaspecialisation
  \\
  Supervisors: & \metropoliainstructors
  \\[6mm]
  \cmidrule[.7pt](l{-.15em}r{5.5em}){1-2}
  \\
  \multicolumn{2}{>{\raggedright}p{15.5cm}}{
    \abstract
  }
  \\
  \\
  Keywords: & \metropoliakeywords
  \\
\end{tabular}
}
\end{otherlanguage}
\restoregeometry
\clearpage


    }{%
    \IfLanguageName{finnish}{%order of abstracts based on main language and spacing hell
        % Abstract in Finnish
% Normally, you should not edit this file. Everything comes from chapter/0info.tex

\pagestyle{empty} %remove page number
\newgeometry{top=1.5cm,left=4cm,right=.5cm}
\begin{otherlanguage}{finnish}
{\large\textbf{Tiivistelmä}}
{\renewcommand{\arraystretch}{1.1}
  \begin{tabular}{@{}p{4.7cm} >{\raggedright\arraybackslash}p{10.8cm}@{}}
  Tekijä: & \makeatletter\@author\makeatother
  \\
  Otsikko: & \otsikko
  \\
  Sivumäärä: & \pageref*{LastPage} sivua + \total{chapter} liitettä
  \\
  Aika: & \pvm
  \\[6.5mm]
  Tutkinto: & \tutkinto
  \\
  Tutkinto-ohjelma: & \kohjelma
  \\
  Ammatillinen pääaine: & \suuntautumis
  \\
  Ohjaajat: & \ohjaajat
  \\[11mm]
  \cmidrule[.7pt](l{-.15em}r{5.5em}){1-2}
  \multicolumn{2}{>{\raggedright}p{15.5cm}}{
  \vspace{.2mm}
  \tiivistelma
  }
  \\
  \\
  Avainsanat: & \avainsanat
  \\
\end{tabular}
}
\end{otherlanguage}
\restoregeometry
\clearpage


        % Abstract in English
% Normally, you should not edit this file. Everything comes from chapter/0info.tex

\pagestyle{empty} %remove page number
\newgeometry{top=1.77cm,left=4cm,right=.5cm}
\begin{otherlanguage}{english}
{\large\textbf{Abstract}}
{\renewcommand{\arraystretch}{1.1}
  \begin{tabular}{@{}p{4.7cm} >{\raggedright\arraybackslash}p{10.8cm}@{}}
  Author: & \makeatletter\@author\makeatother
  \\
  Title: & \makeatletter\@title\makeatother
  \\
  Number of Pages: & \pageref*{LastPage} pages + \total{chapter} appendices
  \\
  Date: & \makeatletter\thedate\makeatother
  \\[6.5mm]
  Degree: & \metropoliadegree
  \\
  Degree Programme: & \metropoliadegreeprogramme
  \\
  Professional Major: & \metropoliaspecialisation
  \\
  Supervisors: & \metropoliainstructors
  \\[9mm]
  \cmidrule[.7pt](l{-.15em}r{5.5em}){1-2}
  \\
  \multicolumn{2}{>{\raggedright}p{15.5cm}}{
    \abstract
  }
  \\
  \\
  Keywords: & \metropoliakeywords
  \\
\end{tabular}
}
\end{otherlanguage}
\restoregeometry
\clearpage


        }{
        % Abstract in English
% Normally, you should not edit this file. Everything comes from chapter/0info.tex

\pagestyle{empty} %remove page number
\newgeometry{top=2.45cm,left=4cm,right=.5cm}
\begin{otherlanguage}{english}
{\large\textbf{Abstract}}
{\renewcommand{\arraystretch}{1.1}
  \begin{tabular}{@{}p{4.7cm} >{\raggedright\arraybackslash}p{10.8cm}@{}}
  Author: & \makeatletter\@author\makeatother
  \\
  Title: & \makeatletter\@title\makeatother
  \\
  Number of Pages: & \pageref*{LastPage} pages + \total{chapter} appendices
  \\
  Date: & \makeatletter\thedate\makeatother
  \\[6.5mm]
  Degree: & \metropoliadegree
  \\
  Degree Programme: & \metropoliadegreeprogramme
  \\
  Professional Major: & \metropoliaspecialisation
  \\
  Supervisors: & \metropoliainstructors
  \\[6mm]
  \cmidrule[.7pt](l{-.15em}r{5.5em}){1-2}
  \\
  \multicolumn{2}{>{\raggedright}p{15.5cm}}{
    \abstract
  }
  \\
  \\
  Keywords: & \metropoliakeywords
  \\
\end{tabular}
}
\end{otherlanguage}
\restoregeometry
\clearpage


        % Abstract in Finnish
% Normally, you should not edit this file. Everything comes from chapter/0info.tex

\pagestyle{empty} %remove page number
\newgeometry{top=2.04cm,left=4cm,right=.5cm}
\begin{otherlanguage}{finnish}
{\large\textbf{Tiivistelmä}}
{\renewcommand{\arraystretch}{1.1}
  \begin{tabular}{@{}p{4.7cm} >{\raggedright\arraybackslash}p{10.8cm}@{}}
  Tekijä: & \makeatletter\@author\makeatother
  \\
  Otsikko: & \otsikko
  \\
  Sivumäärä: & \pageref*{LastPage} sivua + \total{chapter} liitettä
  \\
  Aika: & \pvm
  \\[6.5mm]
  Tutkinto: & \tutkinto
  \\
  Tutkinto-ohjelma: & \kohjelma
  \\
  Ammatillinen pääaine: & \suuntautumis
  \\
  Ohjaajat: & \ohjaajat
  \\[4mm]
  \cmidrule[.7pt](l{-.15em}r{5.5em}){1-2}
  \\
  \multicolumn{2}{>{\raggedright}p{15.5cm}}{
  \tiivistelma
  }
  \\
  \\
  Avainsanat: & \avainsanat
  \\
\end{tabular}
}
\end{otherlanguage}
\restoregeometry
\clearpage


    }
}
%----------------------------------------------------------------------------------------
%	License? Acknowledgement?
%----------------------------------------------------------------------------------------

% Uncomment next line and edit chapters/0license.tex if you want license in your thesis.
%% License of your thesis
% If you wish to explain what it means. When you publish your thesis in https://theseus.fi
% you will be able to choose between some Creative Commons licenses
% https://creativecommons.org
% Adapt this example text to your taste.
% This would also be the right place to explain the license you choose for the code you
% produced for your thesis.

\pagestyle{empty}
\chapter*{Licenses}
\ifdefstring{\thesiscopy}{all}{}{%
  \begin{wrapfigure}{r}{0.3\textwidth}
    \vspace{-20pt}
    \doclicenseImage
  \end{wrapfigure}
 }
\IfLanguageName{finnish}{\copyfi}{\copyen}

That means:

%adapt to the license you choose. In this example: https://creativecommons.org/licenses/by-sa/4.0/deed.en
\textbf{You are free to:}
\begin{itemize}
\item Share \textemdash copy and redistribute the material in any medium or format
\item Adapt \textemdash remix, transform, and build upon the material
\end{itemize}

\textbf{Under the following terms:}
\begin{itemize}
\item \ccAttribution~ Attribution \textemdash You must give appropriate credit, provide a link to the license, and indicate if changes were made. You may do so in any reasonable manner, but not in any way that suggests the licensor endorses you or your use.
\item \ccShareAlike~ ShareAlike \textemdash If you remix, transform, or build upon the material, you must distribute your contributions under the same license as the original.
\item No additional restrictions \textemdash You may not apply legal terms or technological measures that legally restrict others from doing anything the license permits.
\item Any of the above conditions can be waived if you get my permission
\end{itemize}

%Eventually consider few words why you choose such license? E.g. something like:
I decided to publish my thesis work under the Creative Commons Attribution-ShareAlike 4.0 International License because I strongly believe that you as reader deserve the freedom to copy, share and modify this work and if you do modify it, it is fair to give these same permissions to the others. A copy in electronic form of this work can be found in \url{https://some.place} with the \LaTeX{} source.

\clearpage


% Uncomment next line and edit chapters/0acknowledgement.tex if you want acknowledgements.
%% Acknowledgement
% If relevant give a special thanks to the people who supported you during your thesis
% writing.

\pagestyle{empty}
\chapter*{Preface}

A Preface is optional but highly recommended. In this page you may use informal language
such as “This study was a long and hard journey for me but it is done now … finally!”
Remember that after this preface you are not allowed to use informal language and do not use “I” or “You” etc.
in the coming sections (see below).

“I this thesis I studied …” should be something like “This thesis includes a study …”

A Preface, which is typically about half a page long, includes:
-	challenges you faced while writing the thesis
-	things you learned
-	acknowledgments for those who helped you to do your thesis (colleagues, instructor, family, friends, pets ….)

Location, Date
Your Name


\clearpage



%----------------------------------------------------------------------------------------
%	TABLE OF CONTENTS
%----------------------------------------------------------------------------------------

% Table of Content
% Normally, you don't have to modify this file.
% In case spacing of appendix/liite goes crazy, look in ../thesis.tex and look
% for the hack described in APPENDICES.

\makeevenhead{plain}{}{}{}
\makeoddhead{plain}{}{}{}
\pagestyle{empty} %remove page number in toc (if longer than 2 pages)
\tableofcontents*
\pagestyle{empty} %remove page number in toc (if longer than 1 pages)

\clearpage



%list of figure, tables would come here if relevant?

%----------------------------------------------------------------------------------------
%	Lyhenteet / Abbreviation
%----------------------------------------------------------------------------------------

% If you don't use abbreviations/glossary, remove the following line.
% Abbreviation and Glossary
% Normally, you don't have to modify this file. Your abbreviations, etc. goes in
% ../chapters/0abbr.tex file.

\begin{singlespacing}

	% \gsladdall%would add all terms even if not used in your text.
	\addtocontents{toc}{\cftpagenumbersoff{chapter}}
	{
		\vspace{-.46cm}
		\titleformat{\section}
		{\fontsize{13.5pt}{13.5pt}\bfseries\normalfont}
		{\thesection}{.5cm}{}
		\renewcommand*{\glossarypreamble}{\vspace{.7\baselineskip}}
		%Adapt labelwidth (sorry for the ugly hack)
		\setlist[description]{leftmargin=!, labelwidth=4em, before=\doublespacing}
		\IfLanguageName {finnish} {
			\printacronyms[title=Lyhenteet, nogroupskip]
			}{
			\printacronyms[title=List of Abbreviations, nogroupskip]
		}
		\setlist[description]{leftmargin=!, labelwidth=7em, before=\doublespacing}
		\vspace{1cm}
		\printglossary[nogroupskip]
		\setlist[description]{style=standard} % reset settings back to default
	}
	\addtocontents{toc}{\cftpagenumberson{chapter}}
\end{singlespacing}

\clearpage


%----------------------------------------------------------------------------------------
%	CONTENT
%----------------------------------------------------------------------------------------

% Style for the main thesis content.
% Normally, you should not edit this file.

%page number always on top right; also for chapter "title" page
\pagestyle{plain}
\makeevenhead{plain}{}{}{\thepage}
\makeoddhead{plain}{}{}{\thepage}

\setcounter{page}{1} %page 1 should be Introduction

\sloppy % enforce alignment to fully justified


%reset page number to 1, etc.

% Thesis content if you strictly follow the "Final Year Project guide". Of course, you
% can adapt to your specific needs (add more chapter, rename them, etc.).
% Introduction

\chapter{Introduction}

Embedded systems are growing rapidly and becoming really important component of modern life. They
are deployed in devices such as workstations, microprocessors, servers including different consumer
electronics devices, cell phones, interconnected \gls{iot} devices, medical instruments,
autonomous and connected vehicles and many more\cite{camposano1996embedded:ARTICLE}. In a recent study
shows\cite{IoTconne16:online} that the number of such devices are estimated to grow more than
29 billion devices by the year 2030 due to the rapid rise in day-to-day use. Due to this growth, the
amount of related security vulnerabilities are also increasing. The recent Mirai malware attacks,
third party remote exploitation on the automobile, and security attacks on the \gls{ics} are some
examples\cite{muench2018you}.\newline

Software testing with a combination of manual, automation and exploratory method, is a process to
determine whether a software matches it's predefined specification. It is an informative process
when done in the intended environment. Testing helps in uncovering bugs, failures, unwanted errors and
exceptions\cite{819971}. Fuzzing is an automated technique that can be used to generate and input
random malformed data to the system and observe how the software performs\cite{9787842}.
According to \citeauthor{9787842}, due to the closed nature, complexity and high diversity of the
embedded system, fuzzing is not widely used.\newline

The goal for this thesis was to explore the theories of software testing with fuzzing as a testing
technique. This thesis aims to explain what is fuzzing is, its background and its success over the
years discovering bugs. The aim was to do a proof of concept environment using
Google's OSS-Fuzz\cite{GitHubgo49:online} and clusterfuzzlite\cite{ClusterF90:online} project and
evaluate if fuzzing could be implemented for the embedded systems in the case company.\newline

\section{Background}
This study was carried out for a company in Helsinki, Finland which designs and develops embedded secure
silicon \gls{ip} solutions. The solutions range from Key Management Solutions-which helps in securing
key provisioning Management services, Security protocol solutions- to secure the networking services,
\gls{dpa} solutions- to protect side channels attacks and anti-counter fitting products
to \gls{rot} products-highly secure hardware, software and firmware components. \newline


The \gls{rot} solutions which are best described as isolated and separated from the
software/hardware modules that they are meant to defend against the attackers by performing specific critical
functions\cite{Hardware2:online},\cite{RootsofT78:online}. It is the secure foundation and act as a
safeguard against the attackers for \gls{soc},computer, semiconductor and electronic system.
These products are used in securing supply chain solutions, \gls{iot}, cloud and network and
electronics systems.\newline


The company is certified with \gls{iso}- 9001 standard, which is based on the principles of
continual improvement with a strong customer focus and the top management implication and
motivation\cite{ISOISO9048:online}. With a dedicated quality assurance team, infrastructure, the company
is in constant thrive to improve the system testing processes, standards and ultimately
quality and security of the products. The existing system testing includes testing methods such as, manual,
integration, automation, security, exploratory testing. The company follows the \gls{ci} practices,
where the development, integration and testing are done in an incremental way to detect
the errors early, locating the errors faster and testing more. In recent years fuzzing is becoming a
popular and promising technique for the automation security testing. Therefore, the company is
evaluating different fuzzing tools and techniques for its embedded products which will improve the
quality along with the coverage.\newline


\section{Objectives and Purpose}

The objective and aim of this thesis is to evaluate fuzzing techniques, tools as a software automation
technique to improve the quality, security of the products. This study aimed at gathering information
for following questions:\newline

\begin{itemize}
        \item What is fuzzing in the context of software testing? The introduction.
        \item Overview on fuzzing tools and fuzzers available.
        \item Can automation software testing of embedded application improve by fuzzing?
        \item How to build, choose a good fuzzing target? Which fuzzing tool is the best suitable for
        an embedded project?
\end{itemize}


The scope of this thesis was to evaluate fuzzing techniques and give recommendations based on the result by
doing fuzzing on small program. The aim was to get an idea on how it could be implemented in software automation
testing.\newline


For the future integrations and improvements of fuzzing as a test method in the company,
the results of thesis will be used as a baseline.\newline


\section{Thesis Structure And Outline}

The structure and document outline of this thesis are as follows:\newline
Chapter first introduces the subject, theoretical review and benefits of software testing. The
general explanations on test automation, frameworks. The role of software testing in information security,
theories and background of software testing.\newline

Chapter Two explains about fuzzing, why fuzz?, the past and the present, the history. It presents the success of
fuzzing over the years, the vulnerabilities it has discovered, the bugs it has found. It explains about
fuzzing methods, and the implementations.\newline

Chapter three introduces on the fuzzing tools, different types of fuzzers and their comparisons. It
discusses different open source projects and the ongoing research on fuzzing.\newline

Chapter four introduces about the current state analysis of the embedded project. Based on studies of
the previous chapters choosing appropriate fuzzer for the implementations and proof of concept.\newline

Chapter fifth includes the analysis and results of the implementations and case study.\newline

The rest of the chapters include the conclusions, the evaluations and future studies, along with the references
and appendices.\newline

% uncomment what you need.
%% Project Specifications

\chapter{Project Specifications}

Check Final Year Project Guide for the content of Project Specifications chapter.

%% Material and Methods

\chapter{Material and Methods}

Check Final Year Project Guide for the content of Material and Methods chapter.

\clearpage %force the next chapter to start on a new page. Keep that as the last line of your chapter!

%% Theoretical background

\chapter{Theoretical background}

Check Final Year Project Guide for the content of Theoretical background chapter.

\clearpage %force the next chapter to start on a new page. Keep that as the last line of your chapter!

%% Proposed solution
%\clearpage%if the chapter heading starts close to bottom of the page, force a line break and remove the vertical vspace
\vspace{21.5pt}
\chapter{Proposed solution}

Check Final Year Project Guide for the content of Proposed solution chapter.

%% Conclusions
%\clearpage%if the chapter heading starts close to bottom of the page, force a line break and remove the vertical vspace
\vspace{21.5pt}
\chapter{Discussions and Conclusions}
% Discussions
\section{Discussions}

Fuzzing fulfills a dual role in software testing: it is an effective
preliminary screening method for potential vulnerabilities and a sophisticated
tool for probing complex systems. Its effectiveness is evident in the range of
vulnerabilities it uncovers—some are detected quickly, while others may be more
elusive. However, fuzzing alone cannot guarantee system security; the lack of
discovered bugs does not necessarily indicate invulnerability but may suggest
areas for fuzzer enhancement.

The field of fuzzing is continuously evolving, with the introduction of
innovative open-source fuzzers that enhance broader projects. State-of-the-art
tools such as AFL++, libFuzzer, and Atheris reflect a commitment to continuous
improvement and adaptability.

The incorporation of artificial intelligence into fuzzing
strategies~\cite{GoogleOn27:online} introduces a new level of complexity.
AI-powered tools, utilizing machine learning techniques, strive to create
adaptive testing scenarios and improve vulnerability detection rates,
potentially addressing the traditional limitations of fuzzers.

Reflecting on this exploration of fuzzing, it becomes apparent that fuzzing is
not just a component of software testing but a significant field of study.
Overcoming the challenges of understanding fundamental concepts and managing
complex tool configurations has laid a strong foundation for future research,
which will focus on evaluating various fuzzing approaches in different settings.

Embedded systems pose distinct challenges for fuzzing due to their specialized
nature and operational constraints. The variety of hardware platforms requires
the creation of tailored fuzzing tools that can interact with unique hardware
components and communication protocols. Moreover, the constrained computational
resources in embedded devices limit the scope of extensive fuzzing campaigns.

Proprietary operating systems and software stacks in embedded systems, often
optimized for performance and reliability, lack the necessary debugging and
monitoring tools for effective fuzzing. The real-time operations of many
embedded systems also impose strict timing constraints, further complicating the
fuzzing process.

The physical inaccessibility of many embedded devices complicates the
establishment of fuzzing environments and the iterative cycle of test, debug,
and retest, which is crucial for comprehensive vulnerability assessment.
Consequently, there is a significant need for remote and automated fuzzing
solutions tailored to the unique demands of embedded systems.

\subsection*{Advanced Fuzzing Considerations}

A comparative analysis of the aforementioned fuzzing tools indicates distinct
strengths and optimal use cases for each. AFL++ is renowned for its speed and
user-friendliness, libFuzzer's integration with Sanitizers enables thorough code
analysis, and Atheris excels with Python codebases. A deep understanding of each
tool's capabilities is essential for their effective deployment in securing
software.

Case studies demonstrate the practical impact of fuzzing, yet also reveal its
limitations, underscoring the importance of complementary testing methods. These
real-world examples are invaluable for refining fuzzing techniques and
fortifying software systems.

Economic factors play a critical role in the decision to implement fuzzing as a
security strategy. The initial investment in fuzzing infrastructure and
expertise must be balanced against the long-term benefits of early vulnerability
detection and mitigation.

Integrating fuzzing within DevOps, especially in CI/CD pipelines, offers
promising enhancements to software security by embedding automated checks into
the development cycle.


The human element in fuzzing is irreplaceable; the skill and intuition of
security researchers are vital in steering the fuzzing process and interpreting
its outcomes. Advancements in automated fuzzing aim to complement and, in some
cases, surpass human-led testing efforts.


Looking ahead, the field of fuzzing must adapt to emerging technologies like
quantum computing and develop new methodologies to keep pace with the growing
complexity of software architectures.

\section{Conclusions}

This research highlights the significant potential of fuzzing in the realm of
embedded system software testing. Experiments with libFuzzer and AFL++ have
confirmed their ability to swiftly uncover vulnerabilities. These tools, with
their automated reporting features, simplify the complex task of analyzing code
coverage and identifying potential security weaknesses.

The study advocates for the integration of fuzzing into the standard software
development lifecycle, particularly within continuous testing frameworks. As
software intricacy escalates, the necessity for thorough and efficient testing
processes becomes paramount. The strategic implementation of fuzzing tools
ensures that software is not only functional but also secure against potential
threats.

This thesis provides insights into the current landscape of software testing
and establishes a foundation for future testing strategies. The critical need
for evolving testing methodologies is emphasized, and the empirical evidence
presented serves as a basis for adopting advanced testing approaches.

In conclusion, the domain of fuzzing is dynamic and comprehensive. It offers
substantial benefits for software security but also presents numerous challenges
that must be continually addressed. The progressive development of fuzzing tools
and methods, along with a nuanced understanding of their application, remains
crucial for the future of software testing and security.

% Sample content to demonstrate LaTeX command. You will likely delete this line and the
% next \input{sample/*} lines. You are also safe to delete the sample/ folder and its
% content once you refershed your LaTeX skills. Also check the appendix samples.
%sample content to demonstrate LaTeX command.
\vspace{21.5pt}
\chapter{Demonstration Content}\label{demo:content}

This is a chapter to demonstrate some of the \gls{latex} commands that you can use to format your text. If you are new to it, \gls{latex} follows the \gls{wysiwym} idea similar to \gls{html} where you concentrate on the content and structure and leave the formatting and styling to the computer. You will write your content in a plain text file\footnote{which can not get corrupted and can be put under version control} and indicate the structure with commands (similar to \gls{html} tags) then compile it to generate the pdf. Check some books or tutorials such as \LaTeX{} Wikibooks\footnote{\url{https://en.wikibooks.org/wiki/LaTeX}} and try with an online editor such as Overleaf\footnote{\url{https://www.overleaf.com/}} so you do not need to install the compiler on your computer first.

\section{Text, terms and abbreviations}

In \gls{latex}, the \mintinline{tex}{\textbf{bold}} command produces \textbf{bold}, \mintinline{tex}{\textit{italic}}  \textit{italic} and nesting \mintinline{tex}{\textbf{\textit{bolditalic}}} generates \textbf{\textit{bolditalic}}. If the goal is to semantically mark importance, then use \emph{emphasize} with \mintinline{tex}{\emph{emphasize}} command. That would take care of cases such as \mintinline{tex}{\textit{text in italic with \emph{important stuff} inside}} which will compile to \textit{text in italic with \emph{important stuff} inside}.

When one want to use an abbreviation or acronym like \gls{html} using the \mintinline{tex}{\gls{html}} command in \gls{latex}, the first time, it comes with it full version as can be seen in first paragraph of chapter \ref{demo:content} and for all next uses in its short form. Of course, when needed, the full version is available using e.g. the \mintinline{tex}{\acrlong{someID}} command. The defined terms like \gls{maths} use the same \mintinline{tex}{\gls{math}} \gls{latex} command. The Capitalized can be obtained with \mintinline{tex}{\Gls{id}}.

In this template, the abbreviations are defined in the \texttt{chapters/0abbr.tex} file with the \mintinline{tex}{\newacronym{id}{SHORT}{Long Form}} command. There can be many abbreviations and terms defined there, only the ones that are used in the text will be printed in the list of abbreviations (after table of content). And of course, it is the job of the compiler to sort them alphabetically. Should be avoided; but to have all the abbreviations and terms, even the non-used ones, use the \mintinline{tex}{\glsaddall} command before printing the list of abbreviations.

\section{Section with references}
Here is an example how to cite a bibliography entry \cite{kopka:guide} using the \textquotedblleft\textbackslash{}cite\textquotedblright ~\gls{latex} command ~\cite[section 4.2]{tobias:book}. You might also like \textquotedblleft\textbackslash{}citeauthor\textquotedblright and \textquotedblleft\textbackslash{}citedate\textquotedblright ~as demonstrated in figure \ref{fig:latex-cover} caption.  Note that a paragraph is added by forcing a new line.

And let also try the figure (see figure \ref{fig:latex-cover}) and internal reference (with ~\textbackslash{}label and ~\textbackslash{}ref). Alternative text is optained with custom made ~\textbackslash{}AltText command (using pdfcomment and accsupp packages). The reference can be done to any label, for example why not to appendix \ref{appx:first} or to appendix \ref{appx:second}? To note, \gls{latex} will place the figure to the best place (except with forcing). Let them float till the final of final edit\ldots ~then force them to not break a paragraph.%hugly hack... I'm sorry
\begin{figure}[ht]
  \centering
  \AltText{meaningful alternative description (e.g. LaTeX, from typeseting to genrated pdf)}{\includegraphics[width=7.1cm]{LaTeX_cover}}
  \caption{\gls{latex} cover image (Copied from \citeauthor{wikibooks:latex} (\citedate{wikibooks:latex}) \cite{wikibooks:latex}).}
  \label{fig:latex-cover}
\end{figure}

Let's also try a long quote:
From the Universal Declaration of Human Rights:
\begin{quote}
(1) Everyone has the right to education. Education shall be free, at least in the elementary and fundamental stages. Elementary education shall be compulsory. Technical and professional education shall be made generally available and higher education shall be equally accessible to all on the basis of merit.

(2) Education shall be directed to the full development of the human personality and to the strengthening of respect for human rights and fundamental freedoms. It shall promote understanding, tolerance and friendship among all nations, racial or religious groups, and shall further the activities of the United Nations for the maintenance of peace.

(3) Parents have a prior right to choose the kind of education that shall be given to their children. \cite[article 26]{un:udhr}
\end{quote}

\textit{Quisque augue} est, \textbf{elementum ac porttitor} non, porttitor ac orci. Donec hendrerit, ligula ac luctus egestas, sem dolor pretium nunc, sed vehicula magna diam a massa. Donec mattis, arcu et tempor mattis, risus tortor ultrices metus, nec sodales sem dolor eu elit.

\begin{itemize}
  \item Check the guide about the orphan list item: if only first item in the page, force a new page \textbackslash~clearpage before.
  \item When the list items are not sentences, they begin with a lowercase letter, and the last list item ends in a period.
  \item When the list items are sentences, they begin with a capitalized letter, and the list items end in a period. If item of the list contains a long text that spans multiple lines, the left edge aligns automatically.
\end{itemize}

Nullam egestas enim at odio pellentesque bibendum.

\subsection{Subsection}
Donec et sapien ac leo condimentum vulputate id et tellus. Maecenas hendrerit malesuada interdum. Aenean dignissim sem faucibus elit congue faucibus id non risus. Morbi at dui non tortor pellentesque consequat non eget urna. Cras in sapien dui, a tincidunt velit.
\reaction{\label{eq:reaktio}$\underset{\text{+II}}{\ce{2Fe^2+}}$ + $\underset{\text{+I\;-I}}{\ce{H2O2}}$ + $\underset{\text{+I\;-II}}{\ce{2H3O^+}}$ <=> $\underset{\text{+III}}{\ce{2Fe^3+}}$ + $\underset{\text{+I\;-II}}{\ce{4H2O}}$}
Työn aluksi rauta(II)ionit hapetetaan rauta(III)ioneiksi väkevällä vetyperoksidilla, kuten reaktion~\ref{eq:reaktio} hapetusluvuista nähdään (rauta hapettuu, happi pelkistyy).
\reaction{Fe^3+( \emph{aq} ) + 3OH^-( \emph{aq} ) + $(x-1)$H2O( \emph{l} ) -> FeOOH $\cdot$ $x$(H2O)( \emph{s} )}
Rauta(III)ionit saostetaan emäksen (\ce{NH3}) avulla ja saadaan tuotteeksi kidevedellinen rauta(III)hydroksidi. Saatu saostuma pestään \ce{NH4NO3}:lla.
\reaction{FeOOH $\cdot$ $x$(H2O)( \emph{s} ) ->T[$\Delta$900-1000\celsius] Fe2O3( \emph{s} )}

\subsection{Subsection with \texorpdfstring{\Gls{maths}}{Mathematics}}%gls inside chapter/section/... will generate hyperref warning (link inside link in table of content), to avoid that, use the \texorpdfstring
Donec et sapien ac leo condimentum vulputate id et tellus. Maecenas hendrerit malesuada interdum. Aenean dignissim sem faucibus elit congue faucibus id non risus. Morbi at dui non tortor pellentesque consequat non eget urna. Cras in sapien dui, a tincidunt velit. Tertiäärinen butyylikloridi reagoi veden kanssa oheisen reaktion mukaisesti:
\reaction{(CH3)3CCl + 2H2O -> (CH3)3COH+H3O+ +Cl-}
Kyseessä on ensimmäisen kertaluvun reaktio, joten reaktion nopeus on
\begin{align}
v=-\frac{\mathrm{d}[\tn{t-ButCl}]}{\mathrm{d}t}=\frac{\mathrm{d}[\tn{HCl}]}{\mathrm{d}t}=k[\tn{t-ButCl}]
\end{align}
Jos tarkastellaan lähtöaineen t-butyylikloridin häviämistä saadaan
\begin{align}
\frac{\mathrm{d}[\tn{t-ButCl}]}{[\tn{t-ButCl}]}&=-k\mathrm{d}t \\
\int \frac{\mathrm{d}[\tn{t-ButCl}]}{[\tn{t-ButCl}]}&=-k \int \mathrm{d}t \\
\ln \int_{[\tn{t-ButCl}]_0}^{[\tn{t-ButCl}]} [\tn{t-ButCl}]&=-k\int_0^t t \\
\ln \left( \frac{[\tn{t-ButCl}]}{[\tn{t-ButCl}]_0} \right)&=-kt
\end{align}
Ionivahvuus lasketaan kaavalla.
\begin{align}
I&=\frac{1}{2}\cdot\sum z_i^2c_i \\
z_i&= \tn{ionin varausluku} \nonumber \\
c_i&= \tn{ionin konsentraatio} \nonumber
\end{align}
Aktiivisuuskerroin $\gamma_\pm$ lasketaan kaavalla.
\begin{align}
\log \gamma_\pm &= -\left|z_+\cdot z_-\right|A\cdot I^{\frac{1}{2}} \\
A &= \tn{0,509 (lämpötilassa 25\celsius}) \nonumber \\
I &= \tn{ionivahvuus} \nonumber \\
z &= \tn{ionien varaus} \nonumber
\end{align}

\section{Section with Source Code}
Donec et sapien ac leo condimentum vulputate id et tellus. Maecenas hendrerit malesuada interdum. Aenean dignissim sem faucibus elit congue faucibus id non risus. Morbi at dui non tortor pellentesque consequat non eget urna. Cras in sapien dui, a tincidunt velit.

%For sharelatex users, use space instead of tab to avoid ^^I
\begin{code}
  \begin{minted}{php}
<?php
$username = $_POST["username"];
//maybe not?
if ($userName){
  ?>
  <h2>Hello <?= $username ?>!</h2>
  <p>your message got received.</p>
  <?php
}
?>
\end{minted}
\captionof{listing}{Descriptive Caption Text (e.g. this \gls{php} code do blah)}
  \label{code:testphp}
\end{code}


As see in listing \ref{code:testphp}, blah. It is also possible to have code inline, for example \mintinline{sql}{SELECT * FROM user WHERE age >= 18} that was \gls{sql}.
The lisings \ref{code:htmlfull} and \ref{code:htmlpart} show how to load code from an external source file. In the case of listing \ref{code:htmlpart} it only take few line out of the source code file.

\begin{code}
  \inputminted{html}{code/html5_sample.html}
  \captionof{listing}{Some \gls{html} code}
  \label{code:htmlfull}
\end{code}
 Donec et sapien ac leo condimentum vulputate id et tellus. Maecenas hendrerit malesuada interdum. Aenean dignissim sem faucibus elit congue faucibus id non risus.

 \begin{code}
   \inputminted[firstline=3,lastline=6]{html}{code/html5_sample.html}
   \captionof{listing}{The \mintinline{html}{<head>} section of an \gls{html} page}
  \label{code:htmlpart}
\end{code}


 Morbi at dui non tortor pellentesque consequat non eget urna. Cras in sapien dui, a tincidunt velit.


\section{Section with Table}
Donec et sapien ac leo condimentum vulputate id et tellus. Maecenas hendrerit malesuada interdum. Aenean dignissim sem faucibus elit congue faucibus id non risus. Morbi at dui non tortor pellentesque consequat non eget urna. Cras in sapien dui, a tincidunt velit.

\begin{table}[h]
  \centering
  \caption{Some data}%IMPORTANT the caption must be before the tabular, so it will be on top of the table (there are other tricks to force it on top; but this one is straightforward).
  \vspace{-16.5pt}%time to time, spacing between caption and table can go too big...
  \begin{tabular}{| l | >{\centering\arraybackslash}p{.5\textwidth} |}
    \hline
    Test 1 & test 1234 test \\
    \hline
    Some more data comes here & with more values and if the text is very long it will disappear out of the box unless you force the column size :( You can use e.g. \textbackslash raggedright or \textbackslash centering (as in this example) to avoid hyphenization of long words\ldots \\
    \hline
  \end{tabular}
  \label{table:some_data}
\end{table}

As presented in table \ref{table:some_data}: Donec et sapien ac leo condimentum vulputate id et tellus. Maecenas hendrerit malesuada interdum. Aenean dignissim sem faucibus elit congue faucibus id non risus. Morbi at dui non tortor pellentesque consequat non eget urna. Cras in sapien dui, a tincidunt velit.

\begin{table}[h]
  \centering
  \caption{Another table with tabularx}
  \begin{tabularx}{.95\textwidth}{| l | >{\centering\arraybackslash} X |}
    \hline
    Test 1 & test 1234 test \\
    \hline
    Some more data comes here & with more values and if the text is very long it will disappear out of the box unless you force the table size :( \\
    \hline
  \end{tabularx}
  \label{table:some_data2}
\end{table}

As presented in table \ref{table:some_data2}: Donec et sapien ac leo condimentum vulputate id et tellus. Maecenas hendrerit malesuada interdum. Aenean dignissim sem faucibus elit congue faucibus id non risus. Morbi at dui non tortor pellentesque consequat non eget urna. Cras in sapien dui, a tincidunt velit.

Donec et sapien ac leo condimentum vulputate id et tellus. Maecenas hendrerit malesuada interdum. Aenean dignissim sem faucibus elit congue faucibus id non risus. Morbi at dui non tortor pellentesque consequat non eget urna. Cras in sapien dui, a tincidunt velit.

\begin{table}[htbp]
  \centering
  \caption{Booktabs example}
  \vspace{-16.5pt}
    \begin{tabular}{rrrr}
    \toprule
    t (s) & [HCl] & [t-ButCl] & $\ln\frac{[t-ButCl]}{[t-ButCl]_0}$ \\
    \midrule
    0     & 4,02  & 160,88 & 0,00 \\
    10    & 63    & 101,9 & -0,46 \\
    20    & 115,2 & 49,7  & -1,17 \\
    30    & 141,3 & 23,6  & -1,92 \\
    40    & 157,9 & 7     & -3,13 \\
    50    & 161   & 3,9   & -3,72 \\
    60    & 164,3 & 0,6   & -5,59 \\
    70    & 163,5 & 1,4   & -4,74 \\
    80    & 163,8 & 1,1   & -4,99 \\
    90    & 164,1 & 0,8   & -5,30 \\
    100   & 164,3 & 0,6   & -5,59 \\
    \bottomrule
    \end{tabular}
  \label{tab:thisislabel}
\end{table}

As presented in table \ref{tab:thisislabel}: Donec et sapien ac leo condimentum vulputate id et tellus. Maecenas hendrerit malesuada interdum. Aenean dignissim sem faucibus elit congue faucibus id non risus. Morbi at dui non tortor pellentesque consequat non eget urna. Cras in sapien dui, a tincidunt velit.

%Dummy content to use some space...

\chapter{Dummy Content}

Lorem ipsum dolor sit amet, consectetur adipiscing elit. Aliquam aliquam aliquam purus, in ornare nulla imperdiet molestie. Nam tempus erat eu dui rhoncus et vestibulum mi elementum. Ut porttitor elit sit amet justo dignissim sit amet sagittis massa egestas. Mauris sed dolor eget dui fermentum sodales ut eu nibh. 

Quisque augue est, elementum ac porttitor non, porttitor ac orci. Donec hendrerit, ligula ac luctus egestas, sem dolor pretium nunc, sed vehicula magna diam a massa. Donec mattis, arcu et tempor mattis, risus tortor ultrices metus, nec sodales sem dolor eu elit. Nullam egestas enim at odio pellentesque bibendum. 

Donec et sapien ac leo condimentum vulputate id et tellus. Maecenas hendrerit malesuada interdum. Aenean dignissim sem faucibus elit congue faucibus id non risus. Morbi at dui non tortor pellentesque consequat non eget urna. Cras in sapien dui, a tincidunt velit.

\clearpage %force the next chapter to start on a new page. Keep that as the last line of your chapter!

% Sample content to demonstrate tikzpicture

\chapter{Graph}

Lorem ipsum dolor sit amet, consectetur adipiscing elit. Aliquam aliquam aliquam purus, in ornare nulla imperdiet molestie. Nam tempus erat eu dui rhoncus et vestibulum mi elementum. Ut porttitor elit sit amet justo dignissim sit amet sagittis massa egestas. Mauris sed dolor eget dui fermentum sodales ut eu nibh.
\begin{figure}[htbp]
  \centering
    \begin{tikzpicture}
        \pgfplotsset{width=12cm,
        compat=1.3,
        legend style={font=\footnotesize}}
    \begin{axis}[
    xlabel={c (mg/l)},
    ylabel={A},
    legend pos=north west,
    ymajorgrids=true,
    grid style=dashed
]

\addplot [only marks, blue] table {data.dat};
\addplot [no markers, thick, red] table[
x=c,
y={create col/linear regression}] {data.dat};
\addlegendentry{data}
\addlegendentry{%
$\pgfmathprintnumber{\pgfplotstableregressiona}x
\pgfmathprintnumber[print sign]{\pgfplotstableregressionb}$}
\end{axis}
\end{tikzpicture}
\label{fig:stdplot}
\caption{Simple linear regression plot (cannot get $r^2$ value)}
\end{figure}

Quisque augue est, as seen in figure \ref{fig:stdplot} elementum ac porttitor non, porttitor ac orci. Donec hendrerit, ligula ac luctus egestas, sem dolor pretium nunc, sed vehicula magna diam a massa. Donec mattis, arcu et tempor mattis, risus tortor ultrices metus, nec sodales sem dolor eu elit. Nullam egestas enim at odio pellentesque bibendum. 

Donec et sapien ac leo condimentum vulputate id et tellus. Maecenas hendrerit malesuada interdum. Aenean dignissim sem faucibus elit congue faucibus id non risus. Morbi at dui non tortor pellentesque consequat non eget urna. Cras in sapien dui, a tincidunt velit.

\section{Section}

Lorem ipsum dolor sit amet, consectetur adipiscing elit. Aliquam aliquam aliquam purus, in ornare nulla imperdiet molestie. Nam tempus erat eu dui rhoncus et vestibulum mi elementum. 
\tikzstyle{palikka} = [rectangle, rounded corners, minimum width=1cm, minimum height=1cm, text centered, text width=2cm, draw=black, fill=red!30]
\tikzstyle{arrow} = [thick,->,>=stealth]
\begin{figure}[htbp]
\centering
\begin{tikzpicture}[node distance=2.75cm]
\node[label=90:Label] (yksi) [palikka] {Lorem};
\node (kaksi) [palikka, right of=yksi] {ipsum};
\node (kolme) [palikka, below of=kaksi,  yshift=1cm] {dolor};
\node (neljä) [palikka, left of=kolme] {sit};
\node (viisi) [palikka, below of=neljä, yshift=1cm] {amet};
\draw [arrow] (yksi) -- (kaksi);
\draw [arrow] (kaksi) -- node[anchor=west] {tekstiä} (kolme);
\draw [arrow] (kolme) -- (neljä);
\draw [arrow] (neljä) -- (viisi);
\end{tikzpicture}
\caption{Example tikz-picture}
\label{fig:tikz}
\end{figure}

As seen in figure \ref{fig:tikz}, ut porttitor elit sit amet justo dignissim sit amet sagittis massa egestas. Mauris sed dolor eget dui fermentum sodales ut eu nibh. 

\section{Section}

Lorem ipsum dolor sit amet, consectetur adipiscing elit. Aliquam aliquam aliquam purus, in ornare nulla imperdiet molestie. Nam tempus erat eu dui rhoncus et vestibulum mi elementum. Ut porttitor elit sit amet justo dignissim sit amet sagittis massa egestas. Mauris sed dolor eget dui fermentum sodales ut eu nibh. 

\clearpage %force the next chapter to start on a new page. Keep that as the last line of your chapter!


%----------------------------------------------------------------------------------------
%	BIBLIOGRAPHY REFERENCES
%----------------------------------------------------------------------------------------

% Bibliography.
% Normally you do not edit this file.
% To add bibliography in your text, add them first in the biblio.bib file and
% reference them with the \cite{} command in your text. Check also \citeauthor
% and \citeyear (e.g. for copied figure caption), the \cites for multi-cites
% with page numbers,...

\clearpage
%line space
%\singlespacing %removed otherwise the appendix are also single space
\begin{flushleft}
\begin{singlespacing}
\printbibliography[title=\IfLanguageName{finnish}{Lähteet}{References}]
\end{singlespacing}
\end{flushleft}

%for conting the pages
\label{LastPage}~

%avoid that the last page of bib get appendix header
\clearpage



%----------------------------------------------------------------------------------------
%	APPENDICES
%----------------------------------------------------------------------------------------

% Appendix
% Normally, you do not have to edit this file.

%start appendix
\appendix
%no page number for appendix in table of content
\addtocontents{toc}{\cftpagenumbersoff{chapter}}
%appendix sections and subsections not in table of content
\settocdepth{chapter}
%add "Appendices" in the table of content
\addappheadtotoc
%have Appendix 1 (instead of Appendix A)
\renewcommand{\thechapter}{\arabic{chapter}} 

\newcommand\liite[1]{
%each appendix restart page num to one
\setcounter{page}{1}
%special counter for appendix TODO: this is a ugly quick hack :( Should find a better way to count the page per appendix.
\newtotcounter{appx#1}
%overwrite the header
\makeevenhead{plain}{}{}{\appname \thechapter \\ \thepage\,(\stepcounter{appx#1}\total{appx#1})}
\makeoddhead{plain}{}{}{\appname \thechapter \\ \thepage\,(\stepcounter{appx#1}\total{appx#1})}}



%force smaller vertical spacing in table of content
%!!! There can be some fun depending if the appendices have (sub)sections or not :D
% You will have to play with these numbers and eventually add the \vspace line  before
% some \chapter and force another number.
% To add more fun, time to time the table of content get wrong after a build :(
\addtocontents{toc}{\vspace{11pt}}
\pretocmd{\chapter}{\addtocontents{toc}{\protect\vspace{-24pt}}}{}{}

\liite{1}% This is a hack to have right page numbering for each appendix. Make sure to
% use a unique number for each appendix.
% Appendix
% And demonstrate text references and bibliography references in appendix

\chapter{AFL++ Target Program}\label{appx:first}
The following is the detailed implementation of the
target program\cite{FuzzingE54:online}\cite{zephyrsa35:online} used in
the fuzz testing with Zephyr apps by utilizing the AFL++\cite{257204}:

\section*{main.c}
\begin{minted}[linenos,frame=lines,baselinestretch=1.2,breaklines]{c}
#include <stdio.h>
#include <stdbool.h>
#include <stdint.h>
#include <stdlib.h>

/* Fuzz testing is coverage-based, so we want to hide a failure case
 * (a write through a null pointer in this case) down inside a call
 * tree in such a way that it would be very unlikely to be found by
 * randomly-selected input.  But the fuzzer can still find it in
 * linear(-ish) time by discovering each new function along the way
 * and then probing that new space.  The 1 in 2^56 case here would
 * require months-to-years of work for a large datacenter, but the
 * fuzzer gets it in 20 seconds or so. This requires that the code for
 * each case be distinguishable/instrumentable though, which is why we
 * generate the recursive handler functions this way and disable
 * inlining to prevent optimization.
 *
 * https://github.com/zephyrproject-rtos/zephyr/blob/main/samples/subsys/debug/fuzz/src/main.c
 */

int *global_null_ptr;

// crash
static const uint8_t key[] = { 0x9e, 0x21, 0x0c, 0x18, 0x9d, 0xd1, 0x7d };
bool found[sizeof(key)/sizeof(key[0])];

#define LASTKEY (sizeof(key)/sizeof(key[0]) - 1)

#define GEN_CHECK(cur, nxt)                           \
    void check##nxt(uint8_t *data, size_t sz);        \
    void check##cur(uint8_t *data, size_t sz)         \
    {                                                 \
        if (cur < sz && data[cur] == key[cur]) {      \
            if (!found[cur]) {                        \
                printf("Found key %d\n", cur);        \
                found[cur] = true;                    \
            }                                         \
            if (cur == LASTKEY) {                     \
                *global_null_ptr = 0; /* boom! */     \
            } else {                                  \
                check##nxt(data, sz);                 \
            }                                         \
        }                                             \
    }

GEN_CHECK(0, 1)
GEN_CHECK(1, 2)
GEN_CHECK(2, 3)
GEN_CHECK(3, 4)
GEN_CHECK(4, 5)
GEN_CHECK(5, 6)
GEN_CHECK(6, 0)


int main(int argc, char **argv)
{
    /*Matches key
    causes crash*/
    //uint8_t test_data[] = { 0x9e, 0x21, 0x0c, 0x18, 0x9d, 0xd1, 0x7d };
    // no crash
    //uint8_t test_data[] = { 0x9e, 0x21, 0x0c, 0x18, 0x9d, 0xd1, 0x7e };
    // Open the file
   if (argc != 2) {
        printf("Usage: %s <input_file>\n", argv[0]);
        return 1;
    }

    // Open the file
    FILE *file = fopen(argv[1], "rb");
    if (file == NULL) {
        printf("Could not open file %s\n", argv[1]);
        return 1;
    }

    // Determine the size of the file
    fseek(file, 0, SEEK_END);
    size_t size = ftell(file);
    fseek(file, 0, SEEK_SET);

    // Allocate memory for the data
    uint8_t *test_data = malloc(size);
    if (test_data == NULL) {
        printf("Could not allocate memory for test data\n");
        fclose(file);
        return 1;
    }

    // Read the data from the file
    if (fread(test_data, 1, size, file) != size) {
        printf("Error reading file %s\n", argv[1]);
        free(test_data);
        fclose(file);
        return 1;
    }

    // Close the file
    fclose(file);

    // Use the data
    check0(test_data, size);

    // Clean up
    free(test_data);
    return 0;
}

\end{minted}

\section*{CMakeLists.txt}
\begin{minted}{cmake}
    cmake_minimum_required(VERSION 3.20.0)
    project(sample_zephyr_afl)

    option(USE_AFL "Use afl" OFF)

    if(USE_AFL)
        set(CMAKE_C_COMPILER afl-gcc-fast)
        set(CMAKE_C_FLAGS "${CMAKE_C_FLAGS} -g -fpic -fprofile-arcs -ftest-coverage")
    else()
        set(CMAKE_C_COMPILER gcc)
    endif()

    add_executable(zephyr_aflfuzzer src/main.c)
    \end{minted}

\clearpage% Sample content to demonstrate appendix in LaTeX. You
% are safe to delete this lines (and the next samples) once you refreshed your LaTeX
% skills (and safe to delete the sample folder and all its file too).

%\addtocontents{toc}{\vspace{11pt}}%fix vertical space for Table of Content
\liite{2}
% Appendix
% Just to have more appendices and more content

\chapter{Dummy Appendix}\label{appx:second}

Lorem ipsum dolor sit amet, consectetur adipiscing elit. Aliquam aliquam aliquam purus, in ornare nulla imperdiet molestie. Nam tempus erat eu dui rhoncus et vestibulum mi elementum. Ut porttitor elit sit amet justo dignissim sit amet sagittis massa egestas. Mauris sed dolor eget dui fermentum sodales ut eu nibh.

Quisque augue est, elementum ac porttitor non, porttitor ac orci. Donec hendrerit, ligula ac luctus egestas, sem dolor pretium nunc, sed vehicula magna diam a massa. Donec mattis, arcu et tempor mattis, risus tortor ultrices metus, nec sodales sem dolor eu elit. Nullam egestas enim at odio pellentesque bibendum. 

Donec et sapien ac leo condimentum vulputate id et tellus. Maecenas hendrerit malesuada interdum. Aenean dignissim sem faucibus elit congue faucibus id non risus. Morbi at dui non tortor pellentesque consequat non eget urna. Cras in sapien dui, a tincidunt velit.

Lorem ipsum dolor sit amet, consectetur adipiscing elit. Aliquam aliquam aliquam purus, in ornare nulla imperdiet molestie. Nam tempus erat eu dui rhoncus et vestibulum mi elementum. Ut porttitor elit sit amet justo dignissim sit amet sagittis massa egestas. Mauris sed dolor eget dui fermentum sodales ut eu nibh.

Quisque augue est, elementum ac porttitor non, porttitor ac orci. Donec hendrerit, ligula ac luctus egestas, sem dolor pretium nunc, sed vehicula magna diam a massa. Donec mattis, arcu et tempor mattis, risus tortor ultrices metus, nec sodales sem dolor eu elit. Nullam egestas enim at odio pellentesque bibendum. 

Donec et sapien ac leo condimentum vulputate id et tellus. Maecenas hendrerit malesuada interdum. Aenean dignissim sem faucibus elit congue faucibus id non risus. Morbi at dui non tortor pellentesque consequat non eget urna. Cras in sapien dui, a tincidunt velit.

Lorem ipsum dolor sit amet, consectetur adipiscing elit. Aliquam aliquam aliquam purus, in ornare nulla imperdiet molestie. Nam tempus erat eu dui rhoncus et vestibulum mi elementum. Ut porttitor elit sit amet justo dignissim sit amet sagittis massa egestas. Mauris sed dolor eget dui fermentum sodales ut eu nibh.

Quisque augue est, elementum ac porttitor non, porttitor ac orci. Donec hendrerit, ligula ac luctus egestas, sem dolor pretium nunc, sed vehicula magna diam a massa. Donec mattis, arcu et tempor mattis, risus tortor ultrices metus, nec sodales sem dolor eu elit. Nullam egestas enim at odio pellentesque bibendum. 

Donec et sapien ac leo condimentum vulputate id et tellus. Maecenas hendrerit malesuada interdum. Aenean dignissim sem faucibus elit congue faucibus id non risus. Morbi at dui non tortor pellentesque consequat non eget urna. Cras in sapien dui, a tincidunt velit.

Lorem ipsum dolor sit amet, consectetur adipiscing elit. Aliquam aliquam aliquam purus, in ornare nulla imperdiet molestie. Nam tempus erat eu dui rhoncus et vestibulum mi elementum. Ut porttitor elit sit amet justo dignissim sit amet sagittis massa egestas. Mauris sed dolor eget dui fermentum sodales ut eu nibh.

Quisque augue est, elementum ac porttitor non, porttitor ac orci. Donec hendrerit, ligula ac luctus egestas, sem dolor pretium nunc, sed vehicula magna diam a massa. Donec mattis, arcu et tempor mattis, risus tortor ultrices metus, nec sodales sem dolor eu elit. Nullam egestas enim at odio pellentesque bibendum. 

Donec et sapien ac leo condimentum vulputate id et tellus. Maecenas hendrerit malesuada interdum. Aenean dignissim sem faucibus elit congue faucibus id non risus. Morbi at dui non tortor pellentesque consequat non eget urna. Cras in sapien dui, a tincidunt velit.

Lorem ipsum dolor sit amet, consectetur adipiscing elit. Aliquam aliquam aliquam purus, in ornare nulla imperdiet molestie. Nam tempus erat eu dui rhoncus et vestibulum mi elementum. Ut porttitor elit sit amet justo dignissim sit amet sagittis massa egestas. Mauris sed dolor eget dui fermentum sodales ut eu nibh.

Quisque augue est, elementum ac porttitor non, porttitor ac orci. Donec hendrerit, ligula ac luctus egestas, sem dolor pretium nunc, sed vehicula magna diam a massa. Donec mattis, arcu et tempor mattis, risus tortor ultrices metus, nec sodales sem dolor eu elit. Nullam egestas enim at odio pellentesque bibendum. 

Donec et sapien ac leo condimentum vulputate id et tellus. Maecenas hendrerit malesuada interdum. Aenean dignissim sem faucibus elit congue faucibus id non risus. Morbi at dui non tortor pellentesque consequat non eget urna. Cras in sapien dui, a tincidunt velit.

Lorem ipsum dolor sit amet, consectetur adipiscing elit. Aliquam aliquam aliquam purus, in ornare nulla imperdiet molestie. Nam tempus erat eu dui rhoncus et vestibulum mi elementum. Ut porttitor elit sit amet justo dignissim sit amet sagittis massa egestas. Mauris sed dolor eget dui fermentum sodales ut eu nibh.

Quisque augue est, elementum ac porttitor non, porttitor ac orci. Donec hendrerit, ligula ac luctus egestas, sem dolor pretium nunc, sed vehicula magna diam a massa. Donec mattis, arcu et tempor mattis, risus tortor ultrices metus, nec sodales sem dolor eu elit. Nullam egestas enim at odio pellentesque bibendum. 

Donec et sapien ac leo condimentum vulputate id et tellus. Maecenas hendrerit malesuada interdum. Aenean dignissim sem faucibus elit congue faucibus id non risus. Morbi at dui non tortor pellentesque consequat non eget urna. Cras in sapien dui, a tincidunt velit.


\clearpage %force the next appendix to start on a new page. Keep that as the last line of your appendix!


\addtocontents{toc}{\vspace{11pt}}
\liite{3}
% Appendix to demonstrate R integration

\chapter{R examples}

TODO

\clearpage %force the next chapter/appendix to start on a new page. Keep that as the last line of your appendix!



%----------------------------------------------------------------------------------------
%	THIS IS THE END
%----------------------------------------------------------------------------------------
\end{document}
